%%
%% This is file `sigconf.tex',
%% generated with the docstrip utility.
%%
%% The original source files were:
%%
%% samples.dtx  (with options: `all,proceedings,sigconf')
%% 
%% IMPORTANT NOTICE:
%% 
%% For the copyright see the source file.
%% 
%% Any modified versions of this file must be renamed
%% with new filenames distinct from sigconf.tex.
%% 
%% For distribution of the original source see the terms
%% for copying and modification in the file samples.dtx.
%% 
%% This generated file may be distributed as long as the
%% original source files, as listed above, are part of the
%% same distribution. (The sources need not necessarily be
%% in the same archive or directory.)
%%
%%
%% Commands for TeXCount
%TC:macro \cite [option:text,text]
%TC:macro \citep [option:text,text]
%TC:macro \citet [option:text,text]
%TC:envir table 0 1
%TC:envir table* 0 1
%TC:envir tabular [ignore] word
%TC:envir displaymath 0 word
%TC:envir math 0 word
%TC:envir comment 0 0
%%
%% The first command in your LaTeX source must be the \documentclass
%% command.
%%
%% For submission and review of your manuscript please change the
%% command to \documentclass[manuscript, screen, review]{acmart}.
%%
%% When submitting camera ready or to TAPS, please change the command
%% to \documentclass[sigconf]{acmart} or whichever template is required
%% for your publication.
%%
%%
\documentclass[sigconf]{acmart}
%%
%% \BibTeX command to typeset BibTeX logo in the docs
\AtBeginDocument{%
  \providecommand\BibTeX{{%
    Bib\TeX}}}

%% Rights management information.  This information is sent to you
%% when you complete the rights form.  These commands have SAMPLE
%% values in them; it is your responsibility as an author to replace
%% the commands and values with those provided to you when you
%% complete the rights form.
\setcopyright{acmlicensed}
\copyrightyear{2018}
\acmYear{2018}
\acmDOI{XXXXXXX.XXXXXXX}
%% These commands are for a PROCEEDINGS abstract or paper.
\acmConference[Conference acronym 'XX]{Make sure to enter the correct
  conference title from your rights confirmation email}{June 03--05,
  2018}{Woodstock, NY}
%%
%%  Uncomment \acmBooktitle if the title of the proceedings is different
%%  from ``Proceedings of ...''!
%%
%%\acmBooktitle{Woodstock '18: ACM Symposium on Neural Gaze Detection,
%%  June 03--05, 2018, Woodstock, NY}
\acmISBN{978-1-4503-XXXX-X/2018/06}


%%
%% Submission ID.
%% Use this when submitting an article to a sponsored event. You'll
%% receive a unique submission ID from the organizers
%% of the event, and this ID should be used as the parameter to this command.
%%\acmSubmissionID{123-A56-BU3}

%%
%% For managing citations, it is recommended to use bibliography
%% files in BibTeX format.
%%
%% You can then either use BibTeX with the ACM-Reference-Format style,
%% or BibLaTeX with the acmnumeric or acmauthoryear sytles, that include
%% support for advanced citation of software artefact from the
%% biblatex-software package, also separately available on CTAN.
%%
%% Look at the sample-*-biblatex.tex files for templates showcasing
%% the biblatex styles.
%%
%%
%% The majority of ACM publications use numbered citations and
%% references.  The command \citestyle{authoryear} switches to the
%% "author year" style.
%%
%% If you are preparing content for an event
%% sponsored by ACM SIGGRAPH, you must use the "author year" style of
%% citations and references.
%% Uncommenting
%% the next command will enable that style.
%%\citestyle{acmauthoryear}


%%
%% end of the preamble, start of the body of the document source.
\begin{document}

%%
%% The "title" command has an optional parameter,
%% allowing the author to define a "short title" to be used in page headers.
\title{SciVLABench: A Large-Scale Benchmark for Embodied Scientific Agents}

%%
%% The "author" command and its associated commands are used to define
%% the authors and their affiliations.
\author{Author One}
\email{author.one@university.edu}
\affiliation{%
  \institution{University Name}
  \city{City}
  \state{State}
  \country{Country}
}

\author{Author Two}
\correspondingauthor
\email{author.two@university.edu}
\affiliation{%
  \institution{University Name}
  \city{City}
  \state{State}
  \country{Country}
}

%%
%% By default, the full list of authors will be used in the page
%% headers. Often, this list is too long, and will overlap
%% other information printed in the page headers. This command allows
%% the author to define a more concise list
%% of authors' names for this purpose.
\renewcommand{\shortauthors}{Author One et al.}

%%
%% The abstract is a short summary of the work to be presented in the
%% article.
\begin{abstract}
  Embodied agents are increasingly deployed in scientific laboratories, yet systematically benchmarking their capabilities remains an open challenge. While simulation-based evaluation has become the standard paradigm, existing benchmarks suffer from three key limitations. First, they lack a structured task framework, reducing evaluation to isolated operations rather than capturing higher-level competency dimensions. Second, manual task design fundamentally limits scale, failing to cover diverse instructions and complex, operation sequences. Third, the absence of external disturbances---such as human interventions---precludes the assessment of human--robot collaboration.

  To fill this gap, we introduce SciVLABench, a large-scale, comprehensive benchmark for evaluating embodied scientific agents in realistic laboratory settings. In contrast to prior manual efforts, SciVLABench employs an automated task generation pipeline grounded in over 20 curated real-world experimental scenarios. It systematically evaluates eight core capability dimensions---including liquid transfer, physical mixing, thermodynamic control, equipment assembly, and experimental safety---through more than 1,000 complexity-stratified task templates, yielding a task pool 20 times larger than that of the previous largest benchmark. The tasks span basic actions, long-term autonomous experiments, and human--robot collaborative scenarios, providing a fine-grained testbed for model evaluation. SciVLABench is available at \url{https://github.com/kaikai-2022/SciVLABench}.
\end{abstract}

%%
%% The code below is generated by the tool at http://dl.acm.org/ccs.cfm.
%% Please copy and paste the code instead of the example below.
%%
\begin{CCSXML}
<ccs2012>
 <concept>
  <concept_id>00000000.0000000.0000000</concept_id>
  <concept_desc>Do Not Use This Code, Generate the Correct Terms for Your Paper</concept_desc>
  <concept_significance>500</concept_significance>
 </concept>
 <concept>
  <concept_id>00000000.00000000.00000000</concept_id>
  <concept_desc>Do Not Use This Code, Generate the Correct Terms for Your Paper</concept_desc>
  <concept_significance>300</concept_significance>
 </concept>
 <concept>
  <concept_id>00000000.00000000.00000000</concept_id>
  <concept_desc>Do Not Use This Code, Generate the Correct Terms for Your Paper</concept_desc>
  <concept_significance>100</concept_significance>
 </concept>
 <concept>
  <concept_id>00000000.00000000.00000000</concept_id>
  <concept_desc>Do Not Use This Code, Generate the Correct Terms for Your Paper</concept_desc>
  <concept_significance>100</concept_significance>
 </concept>
</ccs2012>
\end{CCSXML}

\ccsdesc[500]{Do Not Use This Code~Generate the Correct Terms for Your Paper}
\ccsdesc[300]{Do Not Use This Code~Generate the Correct Terms for Your Paper}
\ccsdesc{Do Not Use This Code~Generate the Correct Terms for Your Paper}
\ccsdesc[100]{Do Not Use This Code~Generate the Correct Terms for Your Paper}

%% A "teaser" image appears between the author and affiliation
%% information and the body of the document, and typically spans the
%% page.
% \begin{teaserfigure}
%  \includegraphics[width=\textwidth]{sampleteaser}
%   \caption{Seattle Mariners at Spring Training, 2010.}
%   \Description{Enjoying the baseball game from the third-base
%   seats. Ichiro Suzuki preparing to bat.}
%   \label{fig:teaser}
% \end{teaserfigure}

\received{20 February 2007}
\received[revised]{12 March 2009}
\received[accepted]{5 June 2009}

%%
%% This command processes the author and affiliation and title
%% information and builds the first part of the formatted document.
\maketitle


\section{Introduction}
Driven by recent advances in multimodal perception, reasoning, and decision-making, artificial intelligence (AI) is rapidly extending from digital environments into the physical world through embodied agents. While remarkable progress has been achieved in domains such as household assistance and industrial automation, scientific laboratories remain one of the most challenging frontiers for embodied intelligence. Unlike general environments, laboratory tasks require agents to integrate dexterous manipulation, procedural understanding, scientific knowledge, physical reasoning, and safety awareness throughout experimental execution. As embodied agents are increasingly envisioned as autonomous scientific assistants, systematically evaluating their readiness for laboratory environments has become an important and timely research problem.

Developing and evaluating embodied agents in real laboratories is prohibitively expensive, difficult to scale, and inherently risky. High-fidelity simulation therefore provides a natural paradigm for evaluating scientific embodied intelligence in a safe, scalable, and reproducible manner, leading to the emergence of simulation-based laboratory benchmarks. However, building such benchmarks raises a fundamental challenge that is often overlooked: real laboratory protocols are not executable benchmark specifications. They are written for human experimenters and therefore implicitly assume scientific background knowledge, laboratory conventions, physical feasibility, and safety awareness. In contrast, simulation benchmarks require explicit objects, actions, states, constraints, and evaluation criteria that can be executed by embodied agents and verified by the simulator.

The key challenge is thus to bridge this protocol-to-simulation gap. Turning a real experiment into a benchmark task requires more than translating instructions into actions: it requires preserving scientific intent, grounding procedures in embodied interaction, and making experimental outcomes observable and measurable in simulation. For example, a natural-language instruction such as transferring, heating, mixing, or analyzing a sample may correspond to a structured sequence of object selections, parameterized manipulation primitives, state transitions, safety constraints, and success conditions. Without a systematic construction methodology, these implicit structures must be manually rediscovered and reimplemented for every new experiment.

Existing simulation-based laboratory benchmarks are typically constructed through manual task engineering, where each task is individually designed to reproduce a selected laboratory procedure. While this paradigm has enabled valuable progress, it treats every new experiment as an isolated engineering effort. Developers must manually reinterpret the protocol, instantiate the required laboratory assets, implement task-specific logic, and define evaluation rules from scratch. Consequently, benchmark construction remains labor-intensive, difficult to scale, and highly dependent on expert design choices. More fundamentally, current approaches lack a reusable mechanism for transforming real laboratory procedures into executable, verifiable, and evaluable simulation tasks while preserving their scientific semantics.

Realizing such a protocol-driven benchmark construction methodology requires addressing several key challenges. First, laboratory protocols must be semantically represented so that their scientific intent, procedural dependencies, and experimental constraints are preserved throughout benchmark construction. Second, these human-oriented procedures must be grounded into embodied task specifications with explicit objects, actions, states, and success conditions. Third, automatically constructed tasks must be rigorously verified to ensure physical executability, scientific correctness, and reliable evaluation in simulation. Finally, the construction process must scale beyond manually curated task collections, enabling the continuous generation of diverse benchmark tasks from real scientific procedures.

To address these challenges, we present SciVLABench, a protocol-driven benchmark construction framework and large-scale benchmark for evaluating scientific embodied intelligence. Instead of manually engineering each task, SciVLABench formulates laboratory benchmark construction as a protocol-to-task compilation problem. Given natural-language laboratory protocols, the framework progressively transforms scientific procedures into simulation-ready benchmark tasks through semantic grounding, executable task synthesis, and physics-based verification. This compilation process preserves the scientific semantics of real experiments while producing tasks that are executable by embodied agents, verifiable by the simulator, and suitable for standardized evaluation.

Built upon this construction methodology, SciVLABench continuously expands into a comprehensive benchmark covering representative laboratory procedures, diverse scientific operations, and realistic robotic interactions. It currently encompasses more than 20 representative laboratory scenarios derived from contemporary scientific literature and automatically generates over 1,000 task templates covering recurring laboratory operations such as liquid transfer, physical mixing, thermodynamic control, solid material handling and weighing, instrument assembly and spatial organization, anomaly recovery and sensing, chemical analysis, and safety-aware execution. The benchmark further incorporates more than 100 specialized laboratory assets, supports over 30 robotic manipulation primitives, and spans multiple scientific disciplines, including analytical chemistry and biochemistry. Together with automatically generated expert demonstration trajectories for every task template, SciVLABench provides a scalable resource for standardized evaluation and policy learning in scientific laboratory environments.

Our contributions are summarized as follows:

\begin{enumerate}
\item We formulate simulation-based laboratory benchmark construction as a protocol-to-task compilation problem, highlighting the challenge of transforming human-oriented scientific procedures into executable, verifiable, and evaluable embodied tasks.
\item We develop SciVLABench, an automated benchmark construction framework that progressively grounds natural-language laboratory protocols into simulation-ready task specifications through semantic grounding, task synthesis, and physics-based verification.
\item We release a large-scale scientific laboratory benchmark covering representative experimental scenarios, diverse laboratory operations, specialized laboratory assets, robotic manipulation primitives, and automatically generated expert demonstrations, supporting standardized evaluation and policy learning for scientific embodied intelligence.
\end{enumerate}


\section{Related Work}
\label{sec:related-work}
% TODO: Add related work content here

\section{Method}

\subsection{Overview}
\label{sec:overview}

To construct a capability-oriented benchmark, a diverse collection of laboratory task scenarios is required to comprehensively evaluate different scientific capabilities. However, manually designing simulation task templates typically yields only a limited number of isolated scenarios, restricting task diversity and making large-scale benchmark construction difficult. To address this limitation, we develop an automated real-to-simulation task generation pipeline based on LangGraph. Given a natural-language description of a laboratory experiment, the pipeline constructs executable simulation tasks through a closed-loop generation--verification--feedback process, enabling scalable, reproducible, and physically validated benchmark generation.

The proposed pipeline consists of three sequential stages that progressively compile real-world laboratory procedures into executable benchmark tasks. \textbf{Stage I: Environment Construction} transforms a natural-language laboratory protocol into a grounded simulation environment by parsing experimental instructions, grounding referenced entities to laboratory assets, and instantiating a physically consistent laboratory scene. \textbf{Stage II: Task Decomposition and Planning} converts the grounded environment into a verifiable program specification by decomposing the experiment into ordered procedural steps and generating, for each step, both an executable manipulation program and corresponding success conditions. \textbf{Stage III: Simulation Verification and Benchmark Generation} certifies the synthesized specification through physics-grounded execution, condition-guided verification, and iterative refinement, ultimately producing executable task templates, expert demonstrations, and trajectory data. Certified task templates are further diversified through scene randomization before being incorporated into the benchmark.

Methodologically, we formulate automatic benchmark construction as a \textbf{progressive correctness-preserving compilation process}. Instead of directly synthesizing benchmark tasks from natural-language protocols, the pipeline transforms each protocol into three increasingly executable representations: \textbf{a grounded laboratory environment}, \textbf{a verifiable program specification}, and finally \textbf{a certified benchmark instance}. Each stage preserves the semantic correctness established by the previous stage while introducing additional executable structure. Certified task templates are then further diversified through scene randomization before being incorporated into the benchmark, thereby closing the loop from natural-language laboratory protocols to scalable, validated, and reusable benchmark assets.

\subsection{Environment Construction}
\label{sec:environment-construction}

Stage I compiles a natural-language laboratory protocol into a grounded simulation environment. This stage consists of two steps: \textbf{semantic grounding}, which identifies and disambiguates the experimental entities described in the protocol, and \textbf{scene instantiation}, which constructs a physically executable laboratory environment that preserves the semantics of the original procedure.

We first perform semantic grounding to extract laboratory entities and their attributes, including object categories, quantities, materials, and functional roles. Since laboratory protocols often contain multiple instances of the same object category, such as several beakers or test tubes, purely semantic references can become ambiguous during subsequent planning. To avoid this ambiguity, each extracted entity is assigned a globally unique identifier (UID), which is propagated throughout program synthesis, success-condition generation, and benchmark certification. This persistent identifier ensures that all downstream stages refer to the same physical instance.

The grounded entities are then matched against a local laboratory asset repository. If no suitable asset is available, the system retrieves a semantically relevant 3D model from external repositories and converts it into a simulation-ready laboratory object. This conversion includes unifying asset representation, assigning physically meaningful interaction properties such as collision geometry, grasp affordances, articulated structures, and material attributes, and validating geometric consistency, physical plausibility, and manipulation feasibility. As new scenarios are introduced, this retrieval-and-conversion mechanism allows the laboratory asset library to expand beyond manually curated resources.

After all required entities are grounded, the system instantiates the simulation environment according to the spatial relationships specified in the protocol. Object layouts, relative positions, supporting structures, and workspace organization are generated to preserve both procedural semantics and physical feasibility. The resulting scene is not merely a geometric reconstruction of the laboratory setup, but an executable environment in which every object has a consistent semantic identity, valid interaction properties, and a physically plausible spatial configuration. This grounded environment serves as the first executable representation of the laboratory protocol and is passed to the planning stage.

\subsection{Task Decomposition and Planning}
\label{sec:task-decomposition}

Given the grounded environment produced in Stage I, Stage II compiles the experimental procedure into a \textbf{verifiable program specification}. Rather than treating an experiment as a single long-horizon manipulation task, we decompose the procedure into an ordered sequence of semantically complete experimental steps, each corresponding to an independently executable objective. This decomposition reduces long-horizon planning complexity and creates explicit intermediate milestones, enabling process-level evaluation instead of relying solely on final task success.

Based on the step decomposition, the program specification is synthesized through two independent but complementary planning processes. The first process generates an \textbf{execution program} by translating each experimental step into an ordered sequence of robot manipulation primitives selected from a predefined library of more than 30 reusable atomic skills. The planner performs state-aware reasoning over the full experimental procedure, tracking object states and robot interactions across consecutive steps to determine executable manipulation sequences. The execution programs for all steps are then concatenated into a complete expert manipulation trajectory, which can be directly executed in simulation to collect demonstrations for both benchmark evaluation and downstream policy learning.

In parallel, the second planning process generates the \textbf{success conditions} associated with each experimental step. Instead of directly predicting arbitrary verification logic, the planner selects appropriate conditions from a library of expert-designed laboratory verification primitives. These primitives encode experimentally meaningful state constraints, including object configurations, material states, interaction outcomes, and procedural dependencies. The resulting conditions define the expected physical state after each step and provide interpretable criteria for evaluating procedural correctness throughout task execution.

Importantly, the execution program and success conditions are synthesized independently by separate language-model instances. This design naturally decouples task generation from task evaluation, avoiding the undesirable setting in which the same model both generates and validates its own solution. Together, the execution program and success conditions form the program specification of the laboratory procedure. They also enable bidirectional consistency checking during benchmark construction: the execution program verifies whether the success conditions are physically attainable through executable robot behaviors, while the success conditions verify whether the execution program accomplishes the intended experimental objectives. This mutual verification improves the reliability of generated task specifications before they are submitted to simulation certification.

\subsection{Simulation Verification and Benchmark Generation}
\label{sec:simulation-verification}

Stage III addresses a central challenge in automated benchmark construction: \textbf{How can automatically generated tasks be trusted?} Instead of accepting synthesized tasks once they can be executed, SciVLABench requires every generated task to undergo simulation-based certification. In this process, simulation is used not as the final objective, but as a mechanism for verifying that the generated program specification faithfully realizes the intended laboratory procedure under physically executable conditions.

The program specification from Stage II is first combined with the grounded environment from Stage I to form an executable task template. Each template consists of three components: the instantiated simulation environment, the expert execution program, and the corresponding success conditions. The task is then executed in a MuJoCo-based physics simulator, where the robot sequentially performs the synthesized manipulation primitives while success conditions are evaluated throughout the execution process.

Rather than validating only the final task outcome, SciVLABench performs \textbf{topology-aware verification} according to the procedural dependencies of the laboratory protocol. Each success condition is activated only after its prerequisite conditions have been satisfied, ensuring that intermediate experimental states are verified in the correct procedural order. This design enables causal validation of the complete experimental workflow instead of merely checking the final state. A task is certified only when all success conditions are satisfied. During execution, the system records environment configurations, program specifications, execution traces, keyframes, and verification results, providing complete traceability for diagnosis and refinement.

Failed or suspicious generations are not discarded. SciVLABench performs \textbf{LLM-driven secondary inspection} by jointly analyzing execution logs and keyframes captured during condition evaluation. When an inconsistency is identified, the verifier determines the likely failure source and rolls the pipeline back to the responsible compilation stage. Failures caused by unreasonable object placement or invalid spatial layouts trigger environment reconstruction; failures caused by inappropriate manipulation strategies trigger program replanning; and failures arising from unstable physical interactions trigger re-execution or simulation-level refinement. This stage-aware rollback mechanism preserves validated components while correcting only the inconsistent part of the pipeline. The refinement process is repeated for up to three iterations, after which unresolved tasks are flagged for manual inspection.

Finally, certified task templates are transformed into benchmark instances through automatic benchmark generation. To improve policy robustness and support generalization across laboratory environments, certified tasks are diversified by randomizing object poses, spatial arrangements, and material appearances while preserving the semantic validity of the original experimental procedure. The resulting task templates, together with automatically collected expert demonstrations and trajectory data, constitute the final benchmark assets released in SciVLABench. In this way, Stage III completes the correctness-preserving compilation process by promoting synthesized program specifications into certified benchmark instances through verification, diagnosis, and refinement.

\section{Experiment}
% TODO: Add experiment content here

\section{Math Equations}
You may want to display math equations in three distinct styles:
inline, numbered or non-numbered display.  Each of the three are
discussed in the next sections.

\subsection{Inline (In-text) Equations}
A formula that appears in the running text is called an inline or
in-text formula.  It is produced by the \textbf{math} environment,
which can be invoked with the usual
\texttt{{\char'134}begin\,\ldots{\char'134}end} construction or with
the short form \texttt{\$\,\ldots\$}. You can use any of the symbols
and structures, from $\alpha$ to $\omega$, available in
\LaTeX~\cite{Lamport:LaTeX}; this section will simply show a few
examples of in-text equations in context. Notice how this equation:
\begin{math}
  \lim_{n\rightarrow \infty}x=0
\end{math},
set here in in-line math style, looks slightly different when
set in display style.  (See next section).

\subsection{Display Equations}
A numbered display equation---one set off by vertical space from the
text and centered horizontally---is produced by the \textbf{equation}
environment. An unnumbered display equation is produced by the
\textbf{displaymath} environment.

Again, in either environment, you can use any of the symbols and
structures available in \LaTeX\@; this section will just give a couple
of examples of display equations in context.  First, consider the
equation, shown as an inline equation above:
\begin{equation}
  \lim_{n\rightarrow \infty}x=0
\end{equation}
Notice how it is formatted somewhat differently in
the \textbf{displaymath}
environment.  Now, we'll enter an unnumbered equation:
\begin{displaymath}
  \sum_{i=0}^{\infty} x + 1
\end{displaymath}
and follow it with another numbered equation:
\begin{equation}
  \sum_{i=0}^{\infty}x_i=\int_{0}^{\pi+2} f
\end{equation}
just to demonstrate \LaTeX's able handling of numbering.

\section{Figures}

The ``\verb|figure|'' environment should be used for figures. One or
more images can be placed within a figure. If your figure contains
third-party material, you must clearly identify it as such, as shown
in the example below.
\begin{figure}[h]
  \centering
 \includegraphics[width=\linewidth]{sample-franklin}
  \caption{1907 Franklin Model D roadster. Photograph by Harris \&
    Ewing, Inc. [Public domain], via Wikimedia
    Commons. (\url{https://goo.gl/VLCRBB}).}
  \Description{A woman and a girl in white dresses sit in an open car.}
\end{figure}

Your figures should contain a caption which describes the figure to
the reader.

Figure captions are placed {\itshape below} the figure.

Every figure should also have a figure description unless it is purely
decorative. These descriptions convey what’s in the image to someone
who cannot see it. They are also used by search engine crawlers for
indexing images, and when images cannot be loaded.

A figure description must be unformatted plain text less than 2000
characters long (including spaces).  {\bfseries Figure descriptions
  should not repeat the figure caption – their purpose is to capture
  important information that is not already provided in the caption or
  the main text of the paper.} For figures that convey important and
complex new information, a short text description may not be
adequate. More complex alternative descriptions can be placed in an
appendix and referenced in a short figure description. For example,
provide a data table capturing the information in a bar chart, or a
structured list representing a graph.  For additional information
regarding how best to write figure descriptions and why doing this is
so important, please see
\url{https://www.acm.org/publications/taps/describing-figures/}.

\subsection{The ``Teaser Figure''}

A ``teaser figure'' is an image, or set of images in one figure, that
are placed after all author and affiliation information, and before
the body of the article, spanning the page. If you wish to have such a
figure in your article, place the command immediately before the
\verb|\maketitle| command:
\begin{verbatim}
  \begin{teaserfigure}
 \includegraphics[width=\textwidth]{sampleteaser}
    \caption{figure caption}
    \Description{figure description}
  \end{teaserfigure}
\end{verbatim}

\section{Citations and Bibliographies}

The use of \BibTeX\ for the preparation and formatting of one's
references is strongly recommended. Authors' names should be complete
--- use full first names (``Donald E. Knuth'') not initials
(``D. E. Knuth'') --- and the salient identifying features of a
reference should be included: title, year, volume, number, pages,
article DOI, etc.

The bibliography is included in your source document with these two
commands, placed just before the \verb|\end{document}| command:
\begin{verbatim}
  \bibliographystyle{ACM-Reference-Format}
  \bibliography{bibfile}
\end{verbatim}
where ``\verb|bibfile|'' is the name, without the ``\verb|.bib|''
suffix, of the \BibTeX\ file.

Citations and references are numbered by default. A small number of
ACM publications have citations and references formatted in the
``author year'' style; for these exceptions, please include this
command in the {\bfseries preamble} (before the command
``\verb|\begin{document}|'') of your \LaTeX\ source:
\begin{verbatim}
  \citestyle{acmauthoryear}
\end{verbatim}


  Some examples.  A paginated journal article \cite{Abril07}, an
  enumerated journal article \cite{Cohen07}, a reference to an entire
  issue \cite{JCohen96}, a monograph (whole book) \cite{Kosiur01}, a
  monograph/whole book in a series (see 2a in spec. document)
  \cite{Harel79}, a divisible-book such as an anthology or compilation
  \cite{Editor00} followed by the same example, however we only output
  the series if the volume number is given \cite{Editor00a} (so
  Editor00a's series should NOT be present since it has no vol. no.),
  a chapter in a divisible book \cite{Spector90}, a chapter in a
  divisible book in a series \cite{Douglass98}, a multi-volume work as
  book \cite{Knuth97}, a couple of articles in a proceedings (of a
  conference, symposium, workshop for example) (paginated proceedings
  article) \cite{Andler79, Hagerup1993}, a proceedings article with
  all possible elements \cite{Smith10}, an example of an enumerated
  proceedings article \cite{VanGundy07}, an informally published work
  \cite{Harel78}, a couple of preprints \cite{Bornmann2019,
    AnzarootPBM14}, a doctoral dissertation \cite{Clarkson85}, a
  master's thesis: \cite{anisi03}, an online document / world wide web
  resource \cite{Thornburg01, Ablamowicz07, Poker06}, a video game
  (Case 1) \cite{Obama08} and (Case 2) \cite{Novak03} and \cite{Lee05}
  and (Case 3) a patent \cite{JoeScientist001}, work accepted for
  publication \cite{rous08}, 'YYYYb'-test for prolific author
  \cite{SaeediMEJ10} and \cite{SaeediJETC10}. Other cites might
  contain 'duplicate' DOI and URLs (some SIAM articles)
  \cite{Kirschmer:2010:AEI:1958016.1958018}. Boris / Barbara Beeton:
  multi-volume works as books \cite{MR781536} and \cite{MR781537}. A
  presentation~\cite{Reiser2014}. An article under
  review~\cite{Baggett2025}. A
  couple of citations with DOIs:
  \cite{2004:ITE:1009386.1010128,Kirschmer:2010:AEI:1958016.1958018}. Online
  citations: \cite{TUGInstmem, Thornburg01, CTANacmart}.
  Artifacts: \cite{R} and \cite{UMassCitations}.

\section{Acknowledgments}

Identification of funding sources and other support, and thanks to
individuals and groups that assisted in the research and the
preparation of the work should be included in an acknowledgment
section, which is placed just before the reference section in your
document.

This section has a special environment:
\begin{verbatim}
  \begin{acks}
  ...
  \end{acks}
\end{verbatim}
so that the information contained therein can be more easily collected
during the article metadata extraction phase, and to ensure
consistency in the spelling of the section heading.

Authors should not prepare this section as a numbered or unnumbered {\verb|\section|}; please use the ``{\verb|acks|}'' environment.

\section{Appendices}

If your work needs an appendix, add it before the
``\verb|\end{document}|'' command at the conclusion of your source
document.

Start the appendix with the ``\verb|appendix|'' command:
\begin{verbatim}
  \appendix
\end{verbatim}
and note that in the appendix, sections are lettered, not
numbered. This document has two appendices, demonstrating the section
and subsection identification method.

\section{Multi-language papers}

Papers may be written in languages other than English or include
titles, subtitles, keywords and abstracts in different languages (as a
rule, a paper in a language other than English should include an
English title and an English abstract).  Use \verb|language=...| for
every language used in the paper.  The last language indicated is the
main language of the paper.  For example, a French paper with
additional titles and abstracts in English and German may start with
the following command
\begin{verbatim}
\documentclass[sigconf, language=english, language=german,
               language=french]{acmart}
\end{verbatim}

The title, subtitle, keywords and abstract will be typeset in the main
language of the paper.  The commands \verb|\translatedXXX|, \verb|XXX|
begin title, subtitle and keywords, can be used to set these elements
in the other languages.  The environment \verb|translatedabstract| is
used to set the translation of the abstract.  These commands and
environment have a mandatory first argument: the language of the
second argument.  See \verb|sample-sigconf-i13n.tex| file for examples
of their usage.

%%
%% The acknowledgments section is defined using the "acks" environment
%% (and NOT an unnumbered section). This ensures the proper
%% identification of the section in the article metadata, and the
%% consistent spelling of the heading.
\begin{acks}
To Robert, for the bagels and explaining CMYK and color spaces.
\end{acks}


\section*{Ethics and Privacy Statement}

This section of your ACM work should discuss the potential societal
risks that might result from its publication; two to three sentences
related to the findings of your study, or new advancements made
possible by their developed methods. The privacy and ethics statement
should clearly address the broader impacts of their work as it relates
to the authors' interpretation of privacy, fairness, safety, human
rights, data sovereignty, or future misuse and any benefit/risk
trade-off resulting from this research. We acknowledge that some
papers may have minimal societal risks beyond those considered by
institutional review boards, and the dimensions considered by any
review of the user study design or dataset licenses could be provided
in this statement.

%%
%% The next two lines define the bibliography style to be used, and
%% the bibliography file.
\bibliographystyle{ACM-Reference-Format}
\bibliography{sample-base}


%%
%% If your work has an appendix, this is the place to put it.
\appendix

\section{Research Methods}

\subsection{Part One}

Lorem ipsum dolor sit amet, consectetur adipiscing elit. Morbi
malesuada, quam in pulvinar varius, metus nunc fermentum urna, id
sollicitudin purus odio sit amet enim. Aliquam ullamcorper eu ipsum
vel mollis. Curabitur quis dictum nisl. Phasellus vel semper risus, et
lacinia dolor. Integer ultricies commodo sem nec semper.

\subsection{Part Two}

Etiam commodo feugiat nisl pulvinar pellentesque. Etiam auctor sodales
ligula, non varius nibh pulvinar semper. Suspendisse nec lectus non
ipsum convallis congue hendrerit vitae sapien. Donec at laoreet
eros. Vivamus non purus placerat, scelerisque diam eu, cursus
ante. Etiam aliquam tortor auctor efficitur mattis.

\section{Online Resources}

Nam id fermentum dui. Suspendisse sagittis tortor a nulla mollis, in
pulvinar ex pretium. Sed interdum orci quis metus euismod, et sagittis
enim maximus. Vestibulum gravida massa ut felis suscipit
congue. Quisque mattis elit a risus ultrices commodo venenatis eget
dui. Etiam sagittis eleifend elementum.

Nam interdum magna at lectus dignissim, ac dignissim lorem
rhoncus. Maecenas eu arcu ac neque placerat aliquam. Nunc pulvinar
massa et mattis lacinia.

\end{document}
\endinput
%%
%% End of file `sigconf.tex'.
